% Preámbulo modular para presentaciones Beamer (Modelación Estadística)
% Diseño y paleta institucional del Tecnológico de Monterrey

\documentclass[aspectratio=169,xcolor={dvipsnames,table}]{beamer}

\usepackage[utf8]{inputenc}
\usepackage[spanish,mexico]{babel}
\usepackage{amsmath,amssymb,amsfonts,mathtools}
\usepackage{graphicx}
\usepackage{listings}
\usepackage{booktabs}
\usepackage{tikz}
\usetikzlibrary{matrix,arrows,decorations.pathmorphing,shapes.geometric,positioning}

% Paleta Institucional Tecnológico de Monterrey
\definecolor{TecAzul}{HTML}{0039A6}       % Azul TEC: color primario institucional
\definecolor{TecAzulOscuro}{HTML}{1A2E51} % Azul marino oscuro
\definecolor{TecRojo}{HTML}{EC2661}       % Rojo de acento (paleta miTec)
\definecolor{TecGris}{HTML}{646464}       % Gris neutro para comentarios y subtítulos
\definecolor{TecFondoBloque}{HTML}{F4F6F9}% Fondo suave para bloques

% Configuración de Tema Beamer
\usetheme{Boadilla}
\usecolortheme[named=TecAzulOscuro]{structure}
\setbeamercolor{alerted text}{fg=TecRojo}
\setbeamercolor{palette primary}{bg=TecAzulOscuro,fg=white}
\setbeamercolor{palette secondary}{bg=TecAzul,fg=white}
\setbeamercolor{palette tertiary}{bg=TecAzulOscuro!80!black,fg=white}
\setbeamercolor{title}{fg=TecAzulOscuro,bg=TecFondoBloque}
\setbeamercolor{frametitle}{fg=TecAzulOscuro,bg=TecFondoBloque}

% Bloques personalizados elegantes
\setbeamercolor{block title}{bg=TecAzulOscuro,fg=white}
\setbeamercolor{block body}{bg=TecFondoBloque,fg=black}
\setbeamercolor{block title alerted}{bg=TecRojo,fg=white}
\setbeamercolor{block body alerted}{bg=TecRojo!8,fg=black}
\setbeamercolor{block title example}{bg=TecAzul,fg=white}
\setbeamercolor{block body example}{bg=TecAzul!8,fg=black}

% Redondear bordes y añadir sombra sutil
\setbeamertemplate{blocks}[rounded][shadow=true]
\setbeamertemplate{navigation symbols}{}

% Pie de página limpio con número de diapositiva
\setbeamertemplate{footline}{
  \leavevmode%
  \hbox{%
  \begin{beamercolorbox}[wd=.7\paperwidth,ht=2.25ex,dp=1ex,left]{author in head/foot}%
    \hspace*{2ex}\insertshortauthor~~(\insertshortinstitute) \hfill \insertshorttitle
  \end{beamercolorbox}%
  \begin{beamercolorbox}[wd=.3\paperwidth,ht=2.25ex,dp=1ex,right]{date in head/foot}%
    \insertshortdate{}\hspace*{2em}
    \insertframenumber{} / \inserttotalframenumber\hspace*{2ex}
  \end{beamercolorbox}}%
  \vskip0pt%
}

% Configuración de Listings de Python para visualización pedagógica de código
\lstset{
    language=Python,
    basicstyle=\ttfamily\footnotesize,
    keywordstyle=\color{TecAzulOscuro}\bfseries,
    stringstyle=\color{TecRojo},
    commentstyle=\color{TecGris}\itshape,
    showstringspaces=false,
    breaklines=true,
    frame=lines,
    rulecolor=\color{TecAzulOscuro!30},
    backgroundcolor=\color{TecFondoBloque},
    aboveskip=0.5em,
    belowskip=0.5em,
    literate={á}{{\'a}}1 {é}{{\'e}}1 {í}{{\'i}}1 {ó}{{\'o}}1 {ú}{{\'u}}1
             {Á}{{\'A}}1 {É}{{\'E}}1 {Í}{{\'I}}1 {Ó}{{\'O}}1 {Ú}{{\'U}}1
             {ñ}{{\~n}}1 {Ñ}{{\~N}}1 {¿}{{?`}}1 {¡}{{!`}}1
}

% Metadatos institucionales por defecto
\title[Modelación Estadística]{Modelación Estadística para Ciencia de Datos e Ingeniería}
\author[J. Castillo Colmenares]{Juliho David Castillo Colmenares}
\institute[Tec de Monterrey]{Tecnológico de Monterrey \\ Escuela de Ingeniería y Ciencias}
\date{\today}

% Comandos y macros estadísticas compartidas para las presentaciones Beamer (Metropolis)

% Conjuntos numéricos básicos
\newcommand{\R}{\mathbb{R}}
\newcommand{\N}{\mathbb{N}}
\newcommand{\Q}{\mathbb{Q}}
\newcommand{\Z}{\mathbb{Z}}
\newcommand{\set}[1]{\left\{ #1 \right\}}
\newcommand{\abs}[1]{\left|#1\right|}
\newcommand{\norm}[1]{\left\| #1 \right\|}

% Comandos estadísticos y probabilísticos del curso
\newcommand{\Var}{\operatorname{Var}}
\newcommand{\cov}{\operatorname{Cov}}
\newcommand{\corr}{\rho}
\newcommand{\comb}[2]{\begin{pmatrix} #1 \\ #2 \end{pmatrix}}
\newcommand{\s}{\sigma}
\newcommand{\particion}{\mathfrak{P}}
\newcommand{\card}[1]{\#\left( #1 \right)}
\newcommand{\seq}[3]{\set{#1}_{#2}^{#3}}
\newcommand{\E}{\mathbb{E}}
\newcommand{\Prob}{\mathbb{P}}

% Entornos pedagógicos y cajas en español académico natural
\newenvironment{discusion}{%
  \begin{alertblock}{Pregunta de Discusión}%
}{%
  \end{alertblock}%
}

\newenvironment{reflexion}{%
  \begin{alertblock}{Reflexión para el Aula}%
}{%
  \end{alertblock}%
}

\newenvironment{paradoja}{%
  \begin{alertblock}{Paradoja de Exploración}%
}{%
  \end{alertblock}%
}

\newenvironment{motivacion}{%
  \begin{exampleblock}{Motivación: Ciencia de Datos en la Práctica}%
}{%
  \end{exampleblock}%
}

\newenvironment{retoclase}{%
  \begin{block}{Reto Rápido en Equipo}%
}{%
  \end{block}%
}


\title[Estadística Descriptiva]{Sección 1.1: Introducción a la Estadística Descriptiva}
\subtitle{De los Datos Crudos al Conocimiento en Ciencia de Datos}

\begin{document}

% Slide 1: Portada
\begin{frame}
  \titlepage
\end{frame}

% Slide 2: Motivación del Mundo Real
\begin{frame}{¿Por qué necesitamos la Estadística Descriptiva?}
  \begin{motivacion}
    Imagina que eres Científico de Datos en \textbf{Mercado Libre} o \textbf{Netflix}. 
    Cada segundo se generan cientos de miles de clics, compras y tiempos de reproducción.
    
    \vspace{0.3em}
    \textbf{El problema:} Ningún ser humano ni tomador de decisiones puede inspeccionar visualmente una tabla SQL con \textbf{10 millones de filas} para decidir una estrategia.
  \end{motivacion}
  
  \vspace{0.5em}
  La \textbf{Estadística Descriptiva} es nuestra herramienta de compresión de información:
  \begin{itemize}
    \item \textbf{Resumir:} Convierte gigabytes de datos brutos en 3 o 4 métricas interpretables.
    \item \textbf{Detectar:} Saca a la luz patrones ocultos, anomalías graves o fraudes (*outliers*).
    \item \textbf{Comunicar:} Permite reportar hallazgos de forma clara a directivos e ingenieros.
  \end{itemize}
\end{frame}

% Slide 3: Pregunta Gatillo / Discusión en Aula
\begin{frame}{Piénsalo y Discútelo con tu Compañero (2 min)}
  \begin{pregunta}
    Un analista quiere conocer el salario promedio de los ingenieros de software en México. Como no tiene acceso a la nómina de todas las empresas del país, entrevista a \textbf{200 ingenieros} afuera de una torre corporativa en Polanco o San Pedro Garza García.
    \vspace{0.5em}
    \\ \textbf{¿Qué problemas observas en esta metodología?}
    \\ \textbf{¿El resultado representará fielmente la realidad del país?}
  \end{pregunta}
  
  \vspace{0.8em}
  \begin{center}
    \textit{\color{TecAzulOscuro} Reflexión clave: La diferencia entre lo que queremos medir (Población) y lo que realmente podemos observar (Muestra).}
  \end{center}
\end{frame}

% Slide 4: Conceptos Clave - Descriptiva vs Inferencial
\begin{frame}{Dos Ramas Complementarias del Análisis de Datos}
  \begin{columns}[T]
    \begin{column}{0.48\textwidth}
      \begin{block}{Estadística Descriptiva}
        \textbf{¿Qué pasó en estos datos?}
        \begin{itemize}
          \item Se enfoca exclusivamente en \textbf{resumir y organizar} el conjunto de datos observado.
          \item Herramientas: tablas, gráficas de barras, histogramas, promedios, dispersión.
          \item \textbf{Certeza:} 100\% sobre los datos medidos (no hace adivinanzas sobre el futuro).
        \end{itemize}
      \end{block}
    \end{column}
    
    \begin{column}{0.48\textwidth}
      \begin{block}{Estadística Inferencial}
        \textbf{¿Qué podemos generalizar?}
        \begin{itemize}
          \item Utiliza los datos medidos para hacer \textbf{predicciones y generalizaciones} sobre lo que no vimos.
          \item Herramientas: intervalos de confianza, pruebas de hipótesis, regresiones.
          \item \textbf{Incertidumbre:} Cuantifica el margen de error y la probabilidad.
        \end{itemize}
      \end{block}
    \end{column}
  \end{columns}
\end{frame}

% Slide 5: Población vs Muestra & Parámetro vs Estadístico
\begin{frame}{Población vs. Muestra \& Parámetro vs. Estadístico}
  Para no perdernos en la notación matemática, recordemos este puente fundamental:
  
  \vspace{0.4em}
  \begin{center}
    \footnotesize
    \begin{tabular}{p{2.3cm} p{5.5cm} p{5.3cm}}
      \toprule
      \textbf{Concepto} & \textbf{Población (El Universo Total)} & \textbf{Muestra (El Subconjunto Observado)} \\
      \midrule
      \textbf{Definición} & Todos los individuos u objetos de interés & Parte selecta y representativa \\
      \textbf{Medida asociada} & \textbf{Parámetro} (Valor fijo, casi siempre desconocido) & \textbf{Estadístico} (Valor calculado, aleatorio) \\
      \textbf{Media (Promedio)} & $\mu$ (Letra griega mu) & $\bar{X}$ o $\bar{x}$ (X barra) \\
      \textbf{Varianza} & $\sigma^2$ (Sigma al cuadrado) & $S^2$ o $s^2$ (S al cuadrado) \\
      \textbf{Proporción} & $p$ o $\pi$ & $\hat{p}$ (p gorro) \\
      \bottomrule
    \end{tabular}
  \end{center}

  \vspace{0.4em}
  \begin{alertblock}{Regla de oro}
    \textbf{Nunca} calculamos $\mu$ exactamente a menos que hagamos un censo global. ¡Usamos $\bar{X}$ como un telescopio para estimar $\mu$!
  \end{alertblock}
\end{frame}

% Slide 6: Live Coding / Exploración en Python
\begin{frame}[fragile]{Laboratorio en Python: Exploración de Muestras}
  ¿Cómo se traduce esto a código en un flujo real de Ciencia de Datos usando \texttt{numpy} y \texttt{pandas}?
  
  \begin{lstlisting}
import numpy as np
import pandas as pd

# 1. Simulamos una POBLACION oculta de 100,000 usuarios (ej. tiempo en app)
poblacion_mu = 45.0
poblacion_sigma = 12.0
poblacion = np.random.normal(loc=poblacion_mu, scale=poblacion_sigma, size=100000)

# Parametro real (oculto en la practica):
print(f"Parametro poblacional real (mu): {np.mean(poblacion):.2f} min")

# 2. Extraemos una MUESTRA aleatoria de solo 150 usuarios
muestra = np.random.choice(poblacion, size=150, replace=False)

# Estadistico calculado (nuestro estimador):
x_barra = np.mean(muestra)
print(f"Estadistico muestral (X-barra):  {x_barra:.2f} min")
  \end{lstlisting}
\end{frame}

% Slide 7: Quiz / Reto en el Aula
\begin{frame}{Reto Rápido en Equipo: ¿Parámetro o Estadístico?}
  \begin{retoclase}
    Clasifica cada número de los siguientes enunciados como un \textbf{Parámetro ($\mu, \sigma^2, p$)} o un \textbf{Estadístico ($\bar{X}, S^2, \hat{p}$)}:
  \end{retoclase}
  
  \vspace{0.4em}
  \begin{enumerate}
    \item Se escogen al azar 50 transacciones bancarias en un servidor y se encuentra que el monto promedio es de \textbf{\$1,450 MXN}.
    \item El Instituto Nacional de Estadística censó las 35,000 secundarias del país y determinó que el \textbf{12\%} no tiene laboratorio de cómputo.
    \item Un algoritmo de ML entrena con un batch de 512 imágenes y calcula una pérdida (*loss*) promedio de \textbf{0.34}.
    \item Por diseño de ingeniería, una máquina embotelladora llena latas con un volumen promedio teórico exacto de \textbf{355 ml}.
  \end{enumerate}
\end{frame}

% Slide 8: Resumen y Siguientes Pasos
\begin{frame}{Resumen y Conexión con la Sesión 1.2}
  \begin{block}{Lo que debemos recordar hoy (*Takeaways*)}
    \begin{itemize}
      \item La **Estadística Descriptiva** organiza y resume el pasado/presente observado sin incertidumbre.
      \item Siempre debemos distinguir entre el **Parámetro** (la verdad oculta de la Población) y el **Estadístico** (el resumen que obtenemos de nuestra Muestra).
      \item La calidad del análisis depende críticamente de que la muestra sea **representativa** y no sesgada.
    \end{itemize}
  \end{block}

  \vspace{0.5em}
  \begin{exampleblock}{¿Qué veremos en la Sección 1.2?}
    ¿Cuál es el mejor "resumen de un solo número" para nuestros datos?
    \\ Mañana exploraremos las **Medidas de Tendencia Central** (\textit{Media, Mediana y Moda}) y por qué el promedio puede engañarnos drásticamente cuando hay multimillonarios o valores extremos (*outliers*).
  \end{exampleblock}
\end{frame}

\end{document}
